\subsection{Программная система управления и обработки данных}

Программное обеспечение установки организовано в виде трёхуровневой
архитектуры: уровень управления лазерным источником, уровень сбора данных
с АЦП и уровень обработки и визуализации (рис.~\ref{fig:sw-architecture}).
Каждый уровень реализован как самостоятельный программный модуль с
определённым интерфейсом, что обеспечивает модульность и возможность
независимой отладки.

\begin{figure}[H]
    \centering
    \fbox{\parbox[c][0.2\textheight][c]{0.82\textwidth}{\centering
    Место для диаграммы программной архитектуры:\\[4pt]
    Уровень 1: Управление лазерами (Python) $\leftrightarrow$ UART\\
    Уровень 2: Сбор данных (C++) $\leftrightarrow$ USB / TTY\\
    Уровень 3: Обработка и визуализация (Python/PyQtGraph)}}
    \caption{Трёхуровневая программная архитектура установки}
    \label{fig:sw-architecture}
\end{figure}

\subsubsection{Программное обеспечение управления лазерами}

Управление лазерной подсистемой реализовано на языке Python в виде
самостоятельной библиотеки, не зависящей от графического интерфейса. Такая
организация позволяет использовать библиотеку как из графического приложения,
так и из автоматизированных скриптов.

Библиотека обеспечивает подключение к устройству с автоматическим определением
порта и проверкой связи, установку рабочих параметров обоих лазеров
(температура, ток, коэффициенты ПИ-регулятора) с валидацией по допустимым
диапазонам, запуск автоматической перестройки с заданными пределами, шагом и
временными параметрами, а также непрерывный опрос телеметрии с интервалом
5\,мс (температуры, фотодиодные токи, напряжения питания, температура
окружающей среды). Предусмотрена обработка ошибок --- тайм-аутов связи,
ошибок CRC, нарушений валидации, перегрева ТЭО и выхода напряжений за
допуск.

Графический интерфейс отображает временну\'{ю} эволюцию контролируемых
параметров: температуры и их уставки, фотодиодные токи, напряжения шин
питания. Глубина отображения составляет до 100 точек с обновлением в
реальном времени, а общий буфер хранит до 1000 точек для последующего
анализа.

\subsubsection{Программное обеспечение сбора данных}

Сбор данных с модуля E-502 реализован на языке C++ с использованием
библиотеки API производителя (L-Card x502/E-502). Программа сначала
конфигурирует логические каналы АЦП (дифференциальный режим, входной диапазон
$\pm$0{,}2\,В для фазовых измерений, привязка к внешнему тактовому сигналу),
после чего запускает поток данных с аппаратным стробированием. В процессе
работы извлекаются и усредняются окна захвата в реальном времени, результаты
записываются в файлы CSV, а параллельно бинарный поток передаётся через
виртуальный последовательный порт для визуализации.

Дополнительно реализованы две вспомогательные утилиты. Первая ---
\textbf{счётчик тактов} --- измеряет число тактов АЦП между фронтами
стробирующего сигнала, позволяя проверить стабильность тактирования и
корректность синхронизации. Вторая --- \textbf{эмулятор TTY} ---
воспроизводит ранее записанные данные через виртуальный последовательный порт
с регулируемой скоростью, что позволяет разрабатывать и отлаживать
визуализацию без подключения аппаратуры.

\subsubsection{Программное обеспечение обработки и визуализации}

Обработка и визуализация данных реализованы на языке Python с использованием
библиотеки PyQtGraph для построения графиков в реальном времени. Приложение
принимает данные с АЦП, обрабатывает их и одновременно отображает в четырёх
панелях. Панель \textbf{<<текущий свип>>} показывает график частотного отклика
$S_{21}(f_k)$ (модуль или компоненты $I$, $Q$). \textbf{Водопадная диаграмма}
представляет последовательность свипов в виде двумерного изображения
(частота $\times$ время), позволяя отслеживать динамику отклика. Панель
\textbf{<<спектр БПФ>>} отображает результат быстрого преобразования Фурье
текущего свипа с двойной осью абсцисс --- частотой (ГГц) и расстоянием (м),
где ось расстояний вычисляется как
\begin{equation}
    d_n = \frac{n \cdot c}{2\,N_{\mathrm{shift}}\,\Delta f},
    \label{eq:distance-axis}
\end{equation}
причём $c = 299\,792\,458$\,м/с, $N_{\mathrm{shift}}$ --- размер БПФ,
$\Delta f$ --- шаг по частоте. Наконец, панель \textbf{<<B-скан>>} формирует
двумерное радиолокационное изображение путём последовательного накопления
A-сканов (результатов ОБПФ) в кольцевом буфере.

\paragraph{Алгоритм обработки сигнала.}
Обработка каждого принятого свипа проходит несколько последовательных этапов.
Сначала парсер протокола собирает последовательность точек в массив,
интерполируя пропуски средним значением. Затем шаговые индексы пересчитываются
в физические частоты $f\;[\text{ГГц}]$ с помощью полиномиальной калибровки
(см.~раздел~\ref{sec:calibration}). К полученному свипу применяется окно Ханна
для подавления боковых лепестков, после чего вычисляется спектр размерностью
$N = 2048$ точек --- для формирования A-скана выполняется обратное БПФ,
реализующее переход во временну\'{ю} область в соответствии с методикой
предыдущей главы. Из текущего свипа вычитается референсный спектр (измерение в
свободном пространстве) или скользящее медианное среднее по последним 64 свипам.
На последнем этапе амплитуды приводятся к стандартизованному диапазону --- либо
проекцией между калибровочными огибающими, либо поэлементным делением на
калибровочный свип.

\paragraph{Детекция пиков.}
Для автоматического обнаружения отражателей реализован алгоритм поиска
пиков в спектре БПФ. Алгоритм находит наиболее выраженные максимумы
относительно медианной опорной кривой, измеряет их ширину на полувысоте
и пересчитывает положение пика в расстояние. Это позволяет верифицировать
теоретическое разрешение установки при измерениях с известной геометрией
размещения отражателей.
